% SPDX-License-Identifier: CC-BY-SA-4.0
%
% slides.tex — Beamer deck for a 45-minute guest lecture / talk version
% of the course-companion paper.

\documentclass[aspectratio=169,11pt]{beamer}

\usetheme{metropolis}
\usepackage[T1]{fontenc}
\usepackage{lmodern}
\usepackage{listings}
\usepackage{xcolor}

\setbeamertemplate{navigation symbols}{}
\setbeamercovered{transparent}

\lstdefinelanguage{rvasm}{
  morekeywords={addi,add,sub,lw,sw,lbu,beq,bne,blt,jal,jalr,lui,auipc,ecall,fence},
  sensitive=true,
  morecomment=[l]{;},
}

\lstset{
  basicstyle=\ttfamily\scriptsize,
  keywordstyle=\bfseries,
  commentstyle=\color{gray},
  frame=single,
  framerule=0.3pt,
}

\title{Teaching Modern OS Design Through Kernel-Extension VMs}
\subtitle{A MERLIN-V Course Companion}
\author{Weqaar Janjua}
\institute{PacketFive --- University of Limerick}
\date{2026}

\begin{document}

\maketitle

% ----------------------------------------------------------------
\begin{frame}{The pitch}

The single highest-leverage project for a graduate OS course is

\begin{center}
\textbf{a kernel-extension VM.}
\end{center}

\pause
One project covers:
\begin{itemize}
  \item compiler / codegen (the JIT)
  \item program analysis (the verifier)
  \item kernel internals (RCU, anon-inode fds, I-cache, IPI)
  \item UAPI design (the stability promise)
  \item cross-OS portability (the same image on Linux + Zephyr)
\end{itemize}

\pause
No other project hits all five.

\end{frame}

% ----------------------------------------------------------------
\begin{frame}{What's the gap?}

Most OS courses still build a toy kernel:
\begin{itemize}
  \item xv6 / Pintos / OSTEP exercises
  \item Process model, virtual memory, FS, scheduler
\end{itemize}

\pause

The frontier where kernel engineers actually work is
\textbf{controlled extension}:
\begin{itemize}
  \item eBPF (Linux observability + networking)
  \item ChromeOS sandboxing
  \item Cilium, Cilium-derived service meshes
  \item Modern tracing infrastructure
\end{itemize}

\pause

These almost never appear in an OS curriculum at the
``build the verifier'' level.  That's where this course sits.

\end{frame}

% ----------------------------------------------------------------
\begin{frame}{Why not just teach eBPF?}

eBPF has a custom 64-bit RISC-like ISA that exists because in
2014 there was no widely-available simple 64-bit RISC ISA with
production toolchains.

\pause

That precondition no longer holds.

\textbf{RISC-V is:}
\begin{itemize}
  \item Ratified, ABI-stable, multi-vendor
  \item Supported in GCC, LLVM, glibc, musl, Linux, Zephyr
  \item Already the bytecode of many SmartNIC control cores
\end{itemize}

\pause

\textbf{MERLIN-V's premise:} pick RISC-V as the in-kernel
JIT VM's bytecode.  Pedagogically clean.  Toolchain story
collapses to ``stock RISC-V GCC.''

\end{frame}

% ----------------------------------------------------------------
\begin{frame}{Structural consequences}

\begin{itemize}
  \item \textbf{On RISC-V hosts}: the JIT is a pass-through.
        Verified image \emph{is} native machine code.
  \item \textbf{On x86\_64 / arm64}: the JIT still exists, but
        its input is standard RISC-V, not a bespoke bytecode.
  \item \textbf{Toolchain}: no LLVM-BPF backend; no parallel
        gcc-bpf.  Just \texttt{riscv32-unknown-elf-gcc}.
  \item \textbf{Embedded}: the same \texttt{.merlin.o} loads
        into Zephyr on a 400 KiB SRAM MCU.
\end{itemize}

\pause

This is not incremental.  It restructures the
``where does the translation happen'' problem.

\end{frame}

% ----------------------------------------------------------------
\begin{frame}[fragile]{The worked example}

Through the entire course, students return to one tiny program.

\begin{lstlisting}[language=rvasm]
classify:
    lbu  a1, 12(a0)       ; a1 = pkt[12] (EtherType hi byte)
    addi a2, x0, 8        ; a2 = 0x08 (IPv4 high byte)
    beq  a1, a2, .pass
    addi a0, x0, 1        ; return 1 (DROP)
    jalr x0, ra, 0
.pass:
    addi a0, x0, 2        ; return 2 (PASS)
    jalr x0, ra, 0
\end{lstlisting}

\pause

7 instructions.  28 text bytes.

\pause

Every concept in the course shows up in this program:
typed-pointer load, constant materialisation, forward branch,
two return paths, termination.

\end{frame}

% ----------------------------------------------------------------
\begin{frame}{What the verifier must prove}

For this 7-instruction program:

\begin{itemize}
  \item \texttt{a0} on entry is \texttt{PtrCtx(0)}
        $\Rightarrow$ \texttt{lbu a1, 12(a0)} is safe
  \item \texttt{a2 = addi x0, 8} is \texttt{Const(8)}
        $\Rightarrow$ branch is type-checkable
  \item No back-edge $\Rightarrow$ terminates
  \item No \texttt{ecall}, no privileged op
  \item Two return paths, both via \texttt{jalr x0, ra, 0}
\end{itemize}

\pause

\textbf{Pedagogically:} this is where students discover
abstract interpretation.  They concretely instantiate it
themselves in Lab 04 (\textasciitilde 200 lines of C).

\end{frame}

% ----------------------------------------------------------------
\begin{frame}{Course structure: 6 modules, 11 labs}

\begin{enumerate}
  \item \textbf{Setup} (Lab 00)
  \item \textbf{Context} (eBPF dissection $\rightarrow$ RV32 primer)
  \item \textbf{User-space MERLIN-V core} (interpreter $\rightarrow$
        verifier $\rightarrow$ objtool)
  \item \textbf{JITs} (pass-through $\rightarrow$ host JIT to x86\_64)
  \item \emph{Midterm: written design exam, no code}
  \item \textbf{In-kernel runtime} (Linux module $\rightarrow$ Zephyr)
\end{enumerate}

\pause

Plus Lab 10 (hardware bring-up: ESP32-C3 + MPFS Icicle) and
Lab 11 (capstone: student-chosen extension).

\end{frame}

% ----------------------------------------------------------------
\begin{frame}{Lab 03 is the load-bearing one}

Lab 03: \textbf{User-space MERLIN-V interpreter}.

\begin{itemize}
  \item Student starts knowing only C.
  \item Ends with a working RV32I interpreter that runs the
        worked-example program.
  \item Five TODO blocks: ALU-imm, load, store, branch, ecall.
  \item Decoder + dispatch loop + helper trampoline are
        PROVIDED.
\end{itemize}

\pause

\textbf{Why it's load-bearing:} every subsequent lab assumes
this concrete interpreter exists.  The verifier in Lab 04 is
an \emph{abstract} interpreter \textit{over the same dispatch
structure}.  The leap from concrete to abstract is the highest
pedagogical moment in the course.

\end{frame}

% ----------------------------------------------------------------
\begin{frame}{Autograder pipeline}

Every lab is autograded by a 90-line shell script.

\textbf{Locally:}
\begin{itemize}
  \item \texttt{make test} drives the lab's per-test harness
  \item \texttt{runner.sh lab-03} emits one-line JSON
  \item \texttt{score.py} converts JSON to rubric line
\end{itemize}

\pause

\textbf{In CI (GitHub Actions + Azure DevOps):}
\begin{itemize}
  \item Push triggers grading; PR gets a Markdown comment
  \item \textbf{Integrity gate:} canonical instructor files
        overlay the student's tests/Makefile/PROVIDED before
        grading
  \item No way to game the autograder by editing tests
  \item Verified empirically: a fake \texttt{run-all.sh}
        scores 100 without overlay, 14 with overlay.
\end{itemize}

\end{frame}

% ----------------------------------------------------------------
\begin{frame}{Cross-OS portability proof}

The same \texttt{.merlin.o} binary runs on:

\begin{itemize}
  \item \textbf{ESP32-C3-DevKitM-1} (RV32IMC, Zephyr,
        400 KiB SRAM, no MMU)
  \item \textbf{MPFS Icicle Kit} (RV64GC, upstream Linux,
        2 GiB DRAM, kernel module)
\end{itemize}

\pause

\begin{itemize}
  \item Same RV instructions (lbu/addi/beq/jalr encode
        identically on RV32 and RV64).
  \item Different ELF wrappers (Elf32 vs Elf64), same .text.
  \item Same verifier algorithm.
  \item Same helper ABI.
\end{itemize}

\pause

This is the closing demo of the course.  Students see their
own program running on two boards.

\end{frame}

% ----------------------------------------------------------------
\begin{frame}{What this teaches that a toy kernel does not}

\begin{itemize}
  \item \textbf{Stable UAPI} — \texttt{union merlin\_attr} is
        designed to honour the Linux ``never break userspace''
        contract.
  \item \textbf{Verifier soundness as a trust boundary} ---
        bug in JIT corrupts one program; bug in verifier
        compromises the kernel.
  \item \textbf{RCU + refcounting} — real synchronisation
        primitives, used in anger.
  \item \textbf{I-cache coherence} — they discover
        \texttt{flush\_icache\_range} the hard way (a single
        missing flush makes hot-reload silently broken).
  \item \textbf{CI culture} — PR-based, integrity-gated,
        graded.
\end{itemize}

\end{frame}

% ----------------------------------------------------------------
\begin{frame}{Closing}

\begin{block}{Take-aways}
  \begin{itemize}
    \item Kernel-extension VMs are the right OS-curriculum
          vehicle for 2026 onwards.
    \item RISC-V removes a barrier eBPF imposes.
    \item Lab 03 + Lab 04 are the high-leverage moments.
    \item CI + overlay solves academic integrity.
  \end{itemize}
\end{block}

\pause

\textbf{Code, slides, paper:}\\
\url{https://github.com/PacketFive/MERLIN-V/tree/main/docs/academics}

\medskip
CC-BY-SA-4.0 (text), GPL-2.0-only / Apache-2.0 (code).\\
Contributions welcome.

\end{frame}

\end{document}
